%%%%%%%%%%%%%%%%%%%%%%%%%%%%%%%%%%%%%%%%%%%%%%%%%%%%%%%%%%%%%%%%%%%%%%
% Lớp 11 | Hàm số lượng giác và phương trình lượng giác | Các công thức lượng giác
% Dạng 1D1?3-2: Áp dụng công thức cộng
% Số câu: 56  (sinh tự động bởi tach_cau.py; từ nay thầy sửa trực tiếp file này)
%%%%%%%%%%%%%%%%%%%%%%%%%%%%%%%%%%%%%%%%%%%%%%%%%%%%%%%%%%%%%%%%%%%%%%

%%%%%%%%%%%%%%%% TRẮC NGHIỆM 4 PHƯƠNG ÁN %%%%%%%%%%%%%%%%

% nguồn: Ok_Toan11_TracNghiem_Lan2
\begin{ex}%[1D1N3-2]%[ma:95062366]
	Rút gọn biểu thức $M = \tan x - \tan  y$.
	 \choice
		{$M =  \tan(x-y)$}
		{$M =  \dfrac{\sin(x+y)}{\cos x \cos y}$}
		{\True $M = \dfrac{\sin(x-y)}{\cos x  \cos y}$}
		{$M = \dfrac{\tan x -  \tan y}{1+\tan x \tan y}$}
	 \loigiai{
		Ta có $M=\tan x-\tan y=\dfrac{\sin x}{\cos x}-\dfrac{\sin y}{\cos y}=\dfrac{\sin x\cos  y-\cos x\sin y}{\cos x\cos y}=\dfrac{\sin (x-y)}{\cos x\cos y}$.
	} 
\end{ex}

% nguồn: Ok_Toan11_TracNghiem_Lan2
\begin{ex}%[1D1N3-2]%[ma:75ab5fd5]
	Công thức nào sau đây \textbf{sai}? 
	 \choice
		{$\cos (a-b)  = \sin a \sin b + \cos a \cos b$}
		{\True $\cos (a+b) = \sin a  \sin b - \cos a \cos b$}
		{$\sin (a-b)  = \sin a \cos b - \cos a \sin b$}
		{$\sin (a+b) = \sin a \cos b + \cos a \sin b$} 
	\loigiai{ 
		Ta có $\cos  (a+b) = \cos a \cos b - \sin a \sin b$.
	 }
\end{ex}

% nguồn: Ok_Toan11_TracNghiem_Lan2
\begin{ex}%[1D1N3-2]%[ma:bdddc8c1]
	Rút gọn $M = \sin(x-y)\cos y +  \cos(x-y)\sin y$.
	\choice 
		{\True $M = \sin  x$}
		{$M =  \cos x$}
		 {$M = \sin x \cos 2y$}
		{$M = \cos x \cos 2y$} 
	\loigiai{ 
		Áp dụng công thức $\sin (a+b)=\sin a\cos b+\sin  b\cos a$, ta được\\
		$M=\sin (x-y)\cos y+\cos (x-y)\sin y=\sin \left[(x-y)+y \right]=\sin x$. 
	} 
\end{ex}

% nguồn: Ok_TracNghiemLop11_Lan1
\begin{ex}%[1D1N3-2]%[ma:96b840ce]
	Cho $a+b=\dfrac{\pi}{3}$. Giá trị biểu thức $T=\cos  a\cos b-\sin a\sin b$ bằng
	 \choice
		 {$1$}
		 {$-1$}
		 {$\dfrac{\sqrt{3}}{2}$}
		{\True  $\dfrac{1}{2}$}
	 \loigiai{
	Ta có $T=\cos a\cos b-\sin a\sin  b=\cos \left( a+b \right)=\cos \dfrac{\pi }{3}=\dfrac{1}{2}$.}
\end{ex}

% nguồn: Ok_TracNghiemLop11_Lan1
\begin{ex}%[1D1N3-2]%[ma:0ac1c808]
	Công thức nào sau đây {\bf sai}? 
	 \choice
		{\True $\cos  (a+b)=\sin a \sin b-\cos a \cos b$}
		{$\sin (a-b)=\sin a \cos b-\cos a \sin b$} 
		{$\cos (a-b)=\sin a \sin b+\cos a \cos b$} 
		{$\sin (a+b)=\sin a  \cos b+\cos a \sin b$}
	\loigiai{ 
		Ta có  $\cos (a+b)=\cos a \cos b-\sin a \sin b$.
	} 
\end{ex}

% nguồn: Ok_TracNghiemLop11_Lan1
\begin{ex}%[1D1N3-2]%[ma:2e556f05]
	Giá trị lượng giác  $\sin \left(\dfrac{5\pi}{12}\right)$ bằng
	\choice 
		{$0{,}9$} 
		{$\dfrac{\sqrt{2} \left(1  + \sqrt{3}\right)}{2}$}
		{$\dfrac{\sqrt{3}  \left(1 + \sqrt{2}\right)}{4}$}
		{\True  $\dfrac{\sqrt{2} \left(1 + \sqrt{3}\right)}{4}$}
	 \loigiai{
		Ta có 
		\begin{align*}
			\sin \left(\dfrac{5\pi}{12}\right) &= \sin \left(\dfrac{2\pi + 3\pi}{12}\right) = \sin  \left(\dfrac{\pi}{6} + \dfrac{\pi}{4}\right) = \sin \left(\dfrac{\pi}{6}\right) \cos \left(\dfrac{\pi}{4}\right) + \cos \left(\dfrac{\pi}{6}\right) \sin \left(\dfrac{\pi}{4}\right)\\
			&=\dfrac{1}{2}\cdot \dfrac{\sqrt{2}}{2} + \dfrac{\sqrt{3}}{2} \cdot \dfrac{\sqrt{2}}{2} = \dfrac{\sqrt{2}  \left(1 + \sqrt{3}\right)}{4}.
		\end{align*}
	 }
\end{ex}

% nguồn: Ok_TracNghiemLop11_Lan1
\begin{ex}%[1D1N3-2]%[ma:1acc9516]
	Biết $\sin \alpha +\cos \alpha =m$. Tính $P=\cos \left(\alpha  -\dfrac{\pi}{4}\right)$ theo $m$.
	 \choice
		 {$P=2m$}
		{$P=\dfrac{m}{2}$} 
		{\True  $P=\dfrac{m}{\sqrt{2}}$}
		 {$P=m\sqrt{2}$}
	\loigiai{ 
		Ta có $P=\cos \left(\alpha -\dfrac{\pi}{4}\right)=\cos \alpha\cdot \sin \dfrac{\pi}{4}+\sin \alpha  \cdot \cos \dfrac{\pi}{4}=\dfrac{1}{\sqrt{2}}\cos \alpha +\dfrac{1}{\sqrt{2}}\sin \alpha $ .\\
	Do đó $P=\dfrac{1}{\sqrt{2}}\left(\sin \alpha  +\cos \alpha \right)=\dfrac{m}{\sqrt{2}}$.}
\end{ex}

% nguồn: Ok_TracNghiemLop11_Lan1
\begin{ex}%[1D1N3-2]%[ma:e68a8e41]
	Chọn đáp án đúng 
	 \choice
		{\True $\sin a-\sin b=2 \cos \dfrac{a+b}{2} \sin \dfrac{a-b}{2}$} 
		{$\sin a-\sin b=\cos \dfrac{a+b}{2} \sin \dfrac{a-b}{2}$} 
		{$\sin a-\sin b=2  \sin \dfrac{a+b}{2} \cos \dfrac{a-b}{2}$}
		 {$\sin a-\sin b=\sin \dfrac{a+b}{2} \cos \dfrac{a-b}{2}$}
	 \loigiai{
		Công  thức $\sin a-\sin b=2 \cos \dfrac{a+b}{2} \sin \dfrac{a-b}{2}$.
	 }
\end{ex}

% nguồn: Ok_TracNghiemLop11_Lan1
\begin{ex}%[1D1N3-2]%[ma:26db88c0]
	Xét $a$, $b$  là các góc tùy ý, biểu thức $\cos a\cos b+\sin a\sin b$ bằng
	 \choice
		{$\cos (a+b)$} 
		{\True $\cos  (a-b)$}
		{$\sin  (a+b)$}
		 {$\sin (a-b)$}
	 \loigiai{
		Ta có  $\cos a\cos b+\sin a\sin b=\cos (a-b)$.
	} 
\end{ex}

% nguồn: Ok_Toan11_TracNghiem_Lan2
\begin{ex}%[1D1H3-2]%[ma:221d7187]
	Rút  gọn biểu thức $P=\cos \left(120^{\circ}+x\right)+\cos \left(120^{\circ}-x\right)-\cos x$ ta được kết quả là
	 \choice
		 {$0$}
		{$-\cos x$} 
		 {\True $-2\cos x$}
		{$\sin x-\cos  x$}
	 \loigiai{
		Ta có $P=2\cos 120^{\circ} \cos x-\cos x$ $=-\cos x-\cos x=-2\cos  x$.
	 }
\end{ex}

% nguồn: Ok_Toan11_TracNghiem_Lan2
\begin{ex}%[1D1H3-2]%[ma:1bdb7963]
	 Tính $\sin 105^\circ$ ta được
	\choice 
		 {$\dfrac{\sqrt{6}-\sqrt{2}}{4}$}
		 {$-\dfrac{\sqrt{6}-\sqrt{2}}{4}$}
		{\True $\dfrac{\sqrt{6}+\sqrt{2}}{4}$} 
		{$-\dfrac{\sqrt{6}+\sqrt{2}}{4}$} 
	 \loigiai{
		Ta có $\sin 105^\circ = \sin(60^\circ +  45^\circ) = \sin 60^\circ \cos 45^\circ + \cos 60^\circ \sin 45^\circ$.\\
		Suy ra $\sin 105^\circ = \dfrac{\sqrt{3}}{2} \cdot \dfrac{\sqrt{2}}{2} + \dfrac{1}{2} \cdot \dfrac{\sqrt{2}}{2}  = \dfrac{\sqrt{6}+\sqrt{2}}{4}$.
	 }
\end{ex}

% nguồn: Ok_Toan11_TracNghiem_Lan2
\begin{ex}%[1D1H3-2]%[ma:187bf039]
	Giá trị của biểu thức $\cos \dfrac{37\pi}{12}$ bằng 
	 \choice
		 {$\dfrac{\sqrt{6}+\sqrt{2}}{4}$}
		 {$\dfrac{\sqrt{6}-\sqrt{2}}{4}$}
		{\True--$\dfrac{\sqrt{6}+\sqrt{2}}{4}$} 
		{$\dfrac{\sqrt{2}-\sqrt{6}}{4}$} 
	\loigiai{ 
		\begin{eqnarray*}
			\cos \dfrac{37\pi}{12}&=&\cos \left(2\pi+\pi+\dfrac{\pi}{12}  \right)
			=\cos \left(\pi+\dfrac{\pi}{12} \right)\\ 
			 &=&-\cos \left(\dfrac{\pi}{12} \right)=-\cos \left(\dfrac{\pi}{3}-\dfrac{\pi}{4} \right)\\
			&=&-\left(\cos \dfrac{\pi}{3}\cdot \cos \dfrac{\pi}{4}+\sin \dfrac{\pi}{3}\cdot \sin \dfrac{\pi}{4}  \right)\\
			&=&-\dfrac{\sqrt{6}+\sqrt{2}}{4}. 
		\end{eqnarray*}
	 }
\end{ex}

% nguồn: Ok_Toan11_TracNghiem_Lan2
\begin{ex}%[1D1H3-2]%[ma:73b43f9f]
	 Tính $\tan 105^\circ$ ta được
	 \choice
		{\True $-\left(2+\sqrt{3}\right)$} 
		 {$2+\sqrt{3}$}
		{$2-\sqrt{3}$} 
		 {$-\left(2-\sqrt{3}\right)$}
	\loigiai{ 
		Ta có $\tan 105^\circ = \dfrac{\sin 105^\circ}{\cos 105^\circ}= \dfrac{\dfrac{\sqrt{6}+\sqrt{2}}{4}}{-\dfrac{\sqrt{6}-\sqrt{2}}{4}} =  -\dfrac{\sqrt{6}+\sqrt{2}}{\sqrt{6}-\sqrt{2}} = -\left(2+\sqrt{3}\right)$.
	} 
\end{ex}

% nguồn: Ok_Toan11_TracNghiem_Lan2
\begin{ex}%[1D1H3-2]%[ma:c5e1a395]
	Có bao  nhiêu đẳng thức dưới đây là đồng nhất thức?
	\begin{itemize}
		\begin{multicols}{2}
			\item $\cos x - \sin x =  \sqrt{2} \sin \left(x + \dfrac{\pi}{4}\right)$
			\item $\cos x - \sin x = \sqrt{2}  \cos \left(x + \dfrac{\pi}{4}\right)$
			\item  $\cos x - \sin x = \sqrt{2} \sin \left(x - \dfrac{\pi}{4}\right)$
			 \item $\cos x - \sin x = \sqrt{2} \sin \left(\dfrac{\pi}{4} - x\right)$
		\end{multicols}
	\end{itemize}
	\choice 
		 {$1$}
		 {\True $2$}
		{$3$} 
		 {$4$}
	\loigiai{ 
		Ta có  $\cos x - \sin x = \sqrt{2} \cos \left(x + \dfrac{\pi}{4}\right)=\sqrt{2}\cos\left[\dfrac{\pi}{2}-\left(\dfrac{\pi}{4}-x\right)\right]=\sqrt{2}\sin\left(\dfrac{\pi}{4}-x\right)$.
	 }
\end{ex}

% nguồn: Ok_Toan11_TracNghiem_Lan2
\begin{ex}%[1D1H3-2]%[ma:f87e5cc3]
	Cho hai góc nhọn $a$ và $b$. Biết $\cos a=\dfrac{1}{3}$,  $\cos b=\dfrac{1}{4}$.
	Giá  trị $\cos \left(a+b\right).\cos \left(a-b\right)$ bằng
	\choice 
		{$-\dfrac{113}{144}$} 
		{$-\dfrac{115}{144}$} 
		{$-\dfrac{117}{144}$} 
		{\True $-\dfrac{119}{144}$} 
	 \loigiai{
		\begin{eqnarray*}
			&&  \cos \left(a+b\right)\cdot \cos \left(a-b\right)\\
			 &=&\dfrac{1}{2} \left(\cos 2a+\cos 2b\right)\\
			&=&\cos ^2a+\cos  ^2b-1\\
			&=&\left(\dfrac{1}{3} \right)^2+\left(\dfrac{1}{4}  \right)^2-1\\
			&=&-\dfrac{119}{144}. 
		\end{eqnarray*}
	 }
\end{ex}

% nguồn: Ok_Toan11_TracNghiem_Lan2
\begin{ex}%[1D1H3-2]%[ma:9a3b9f3a]
	Nếu $\tan x=0{,}5;$ $\sin y=\dfrac{3}{5}\, \left(0^\circ< y < 90^{\circ}\right)$ thì  $\tan \left(x+y\right)$ bằng
	 \choice
		{\True  $2$}
		{$3$} 
		 {$4$}
		{$5$} 
	\loigiai{ 
		Do $0^\circ<  y < 90^{\circ}$ nên $\cos y>0$.\\
		$ \Rightarrow \cos y= \sqrt{1-\sin^2 y  }= \sqrt{1- \dfrac{9}{25}}= \dfrac{4}{5} \Rightarrow \tan y=\dfrac{3}{4}.$
		\\ 
		$\tan \left(x+y\right)=\dfrac{\tan x+\tan y}{1-\tan x\cdot \tan y}=\dfrac{\dfrac{1}{2}+\dfrac{3}{4}}{1-\dfrac{1}{2}\cdot  \dfrac{3}{4}}=2.$
	} 
\end{ex}

% nguồn: Ok_Toan11_TracNghiem_Lan2
\begin{ex}%[1D1H3-2]%[ma:993d3449]
	Giá  trị đúng của $A=\cos \dfrac{2\pi}{7}+\cos \dfrac{4\pi}{7}+\cos \dfrac{6\pi}{7}$ bằng
	\choice 
		{$\dfrac{1}{2}$} 
		 {\True $-\dfrac{1}{2}$}
		 {$\dfrac{1}{4}$}
		 {--$\dfrac{1}{4}$}
	\loigiai{ 
		Ta có \begin{eqnarray*}
		2\sin \dfrac{\pi}{7}\cdot A & = & 2\sin \dfrac{\pi}{7}\cdot \cos \dfrac{2\pi}{7}+  2\sin \dfrac{\pi}{7}\cdot \cos \dfrac{4\pi}{7}+ 2\sin \dfrac{\pi}{7}\cdot\cos \dfrac{6\pi}{7}\\
		&= & \sin \dfrac{3 \pi}{7}- \sin  \dfrac{\pi }{7} + \sin \dfrac{5\pi}{7} - \sin\dfrac{3\pi}{7}+ \sin \pi - \sin \dfrac{5\pi}{7} \\
		& = & -  \sin \dfrac{\pi}{7}.
	\end{eqnarray*}
	\noindent Vậy  $A= -\dfrac{1}{2}$.
} 
\end{ex}

% nguồn: Ok_Toan11_TracNghiem_Lan2
\begin{ex}%[1D1H3-2]%[ma:1a5539f5]
	Cho hai góc  $\alpha$, $\beta$ thỏa mãn $\sin \alpha=\dfrac{5}{13}$, $\left(\dfrac{\pi}{2}<\alpha<\pi\right)$ và $\cos \beta=\dfrac{3}{5}$, $\left(0<\beta<\dfrac{\pi}{2}\right)$. Tính giá trị đúng của $\cos (\alpha-\beta)$.
	 \choice
		 {$\dfrac{16}{65}$}
		{$-\dfrac{18}{65}$} 
		{$\dfrac{18}{65}$} 
		{\True $-\dfrac{16}{65}$} 
	\loigiai{ 
		Ta có $\sin \alpha=\dfrac{5}{13}$, $\left(\dfrac{\pi}{2}<\alpha<\pi\right)$ nên $\cos  \alpha=-\sqrt{1-\left(\dfrac{5}{13}\right)^{2}}=-\dfrac{12}{13}$.\\
		Ta có 
		$\cos \beta=\dfrac{3}{5}$, $\left(0<\beta<\dfrac{\pi}{2}\right)$  nên $\sin \beta=\sqrt{1-\left(\dfrac{3}{5}\right)^{2}}=\dfrac{4}{5}$.\\
	Vậy $\cos (\alpha-\beta)=\cos \alpha \cos \beta+\sin \alpha \sin  \beta=-\dfrac{12}{13} \cdot \dfrac{3}{5}+\dfrac{5}{13} \cdot \dfrac{4}{5}=-\dfrac{16}{65}$.}
\end{ex}

% nguồn: Ok_Toan11_TracNghiem_Lan2
\begin{ex}%[1D1H3-2]%[ma:2dde0ab3]
	Giá trị của biểu  thức $\cos \dfrac{37\pi}{12}$ bằng
	\choice 
		 {$\dfrac{\sqrt{6}+\sqrt{2}}{4}$}
		{$\dfrac{\sqrt{6}-\sqrt{2}}{4}$} 
		 {\True $-\dfrac{\sqrt{6}+\sqrt{2}}{4}$}
		 {$\dfrac{\sqrt{2}-\sqrt{6}}{4}$}
	\loigiai{ 
		Ta có $\cos \dfrac{37\pi}{12} = \cos \left(\dfrac{\pi}{2} + \dfrac{7\pi}{12}+2\pi\right) = \cos \left(\dfrac{\pi}{2} + \dfrac{7\pi}{12}\right)= \cos \dfrac{\pi}{2} \cos \dfrac{7\pi}{12} - \sin \dfrac{\pi}{2}  \sin \dfrac{7\pi}{12} = -\dfrac{\sqrt{6}+\sqrt{2}}{4}$.
	 }
\end{ex}

% nguồn: Ok_Toan11_TracNghiem_Lan2
\begin{ex}%[1D1H3-2]%[ma:6ba2d2c5]
	Rút gọn $M =  \cos(a+b)\cos(a-b) + \sin(a+b)\sin(a-b)$.
	\choice 
		{\True  $M=1-2\sin^2 b$}
		{$M=1+2\sin^2  b$}
		{$M=\cos 4b$} 
		{$M=\sin 4b$} 
	\loigiai{ 
		Áp dụng công thức $\cos x\cos  y+\sin x\sin y=\cos (x-y)$, ta được\\
		$M=\cos  (a+b)\cos (a-b)+\sin (a+b)\sin (a-b)=\cos (a+b-(a-b))=\cos 2b=1-2\sin^2 b$.
	} 
\end{ex}

% nguồn: Ok_Toan11_TracNghiem_Lan2
\begin{ex}%[1D1H3-2]%[ma:a9d140ad]
	Rút gọn biểu thức $\cos 54^\circ  \cos 4^\circ - \cos 36^\circ \cos 86^\circ$, ta được
	\choice 
		 {$\cos 50^\circ$}
		{\True  $\cos 58^\circ$}
		{$\sin 50^\circ$} 
		 {$\sin 58^\circ$}
	 \loigiai{
		Ta có $\cos 54^\circ \cos 4^\circ - \cos 36^\circ \cos  86^\circ = \cos 54^\circ \cos 4^\circ - \sin 54^\circ \sin 4^\circ= \cos (54^\circ + 4^\circ) = \cos 58^\circ$.
	 }
\end{ex}

% nguồn: Ok_Toan11_TracNghiem_Lan2
\begin{ex}%[1D1H3-2]%[ma:ab1c9c54]
	Biểu thức $M=\cos \left(-53^{\circ} \right)\cdot \sin \left(-337^{\circ} \right)+\sin  \left(307^{\circ} \right)\cdot \sin \left(113^{\circ} \right)$ có giá trị bằng
	 \choice
		{\True $-\dfrac{1}{2}$} 
		 {$\dfrac{1}{2}$}
		 {$-\dfrac{\sqrt{3}}{2}$}
		{$\dfrac{\sqrt{3}}{2}$} 
	\loigiai{ 
		 Ta có
		\begin{eqnarray*}
			M &= &\cos \left(-53^{\circ}  \right)\cdot \sin \left(-337^{\circ} \right)+\sin \left(307^{\circ} \right)\cdot \sin \left(113^{\circ} \right) \\
			& = & \cos 53^\circ \cdot \sin 23^\circ  + \sin (-53 ^\circ) \cdot \sin 67^\circ\\
			& = & \sin 23 ^\circ \cdot \cos 53^\circ -  \sin 53^\circ \cdot \cos 23^\circ \\
			& = &  -\sin (53^\circ - 23^\circ)\\
			& = & -  \sin 30^\circ =- \dfrac{1}{2}
			 .
		\end{eqnarray*}
	} 
\end{ex}

% nguồn: Ok_Toan11_TracNghiem_Lan2
\begin{ex}%[1D1H3-2]%[ma:d3be0775]
	 Kết quả nào sau đây \textbf{sai}?
	\choice 
		 {$\sin x+\cos x=\sqrt{2} \sin \left(x+\dfrac{\pi}{4}\right)$}
		 {$\sin x-\cos x=-\sqrt{2} \cos \left(x+\dfrac{\pi}{4}\right)$}
		{\True $\sin 2 x+\cos 2 x=\sqrt{2}  \sin \left(2 x-\dfrac{\pi}{4}\right)$}
		 {$\sin 2 x+\cos 2 x=\sqrt{2} \cos \left(2 x-\dfrac{\pi}{4}\right)$}
	 \loigiai{
		 Ta có
		\begin{align*}
			\sin 2 x+\cos 2 x&=\sqrt{2}\left(\dfrac{1}{\sqrt{2}} \sin  2 x+\dfrac{1}{\sqrt{2}} \cos 2 x\right)\\
			&=\sqrt{2}\left(\cos \dfrac{\pi}{4} \sin  2 x+\sin \dfrac{\pi}{4} \cos 2 x\right)\\
			&=\sqrt{2} \sin \left(2 x+\dfrac{\pi}{4}\right) \neq \sqrt{2} \sin \left(2  x-\dfrac{\pi}{4}\right).
	\end{align*}}
\end{ex}

% nguồn: Ok_Toan11_TracNghiem_Lan2
\begin{ex}%[1D1H3-2]%[ma:aae4957e]
	Cho $M=\tan x-\tan  y$, chọn khẳng định đúng.
	 \choice
		{$M=\tan \left(x-y\right)$} 
		 {$M=\dfrac{\sin \left(x+y\right)}{\cos x\cdot \cos y}$}
		{\True $M=\dfrac{\sin \left(x-y\right)}{\cos x\cdot \cos  y}$}
		{$M=\dfrac{\tan x-\tan y}{1+\tan  x\cdot \tan y}$}
	\loigiai{ 
		Ta có $M=\tan  x-\tan y=\dfrac{\sin x}{\cos x}-\dfrac{\sin y}{\cos y}=\dfrac{\sin x\cos y-\cos x\sin y}{\cos x\cos y}$ $=\dfrac{\sin \left(x-y\right)}{\cos x\cos y}$.
	} 
\end{ex}

% nguồn: Ok_Toan11_TracNghiem_Lan2
\begin{ex}%[1D1H3-2]%[ma:440a5679]
	Rút gọn biểu thức: $\sin \left(a-17^{\circ} \right)\cdot \cos  \left(a+13^{\circ} \right)-\sin \left(a+13^{\circ} \right)\cdot \cos \left(a-17^{\circ} \right)$,
	ta được kết  quả nào sau đây?
	\choice 
		{$\sin 2a$} 
		{$\cos  2a$}
		{\True  $-\dfrac{1}{2}$}
		 {$\dfrac{1}{2}$}
	 \loigiai{
		Ta  có
		\begin{eqnarray*}
			\sin \left(a-17^{\circ} \right)\cdot \cos \left(a+13^{\circ} \right)-\sin  \left(a+13^{\circ} \right)\cdot \cos \left(a-17^{\circ} \right) &=& \sin \left[\left(a-17^{\circ} \right)-\left(a+13^{\circ} \right)\right]\\
			&=&\sin \left(-30^{\circ} \right)=-\dfrac{1}{2}. 
		\end{eqnarray*}
	} 
\end{ex}

% nguồn: Ok_Toan11_TracNghiem_Lan2
\begin{ex}%[1D1H3-2]%[ma:344c216d]
	Giả sử $A=\tan x\cdot \tan\left(\dfrac{\pi}{3}-x\right)\cdot \tan\left(\dfrac{\pi}{3}+x\right)$ được rút gọn  thành $A=\tan nx$ khi đó $n$ bằng
	\choice 
		 {$ 2$}
		{$ 1$} 
		{$ 4$} 
		{\True $  3$}
	 \loigiai{
		 Ta có
		\allowdisplaybreaks\begin{eqnarray*}
		A&=&\tan x\cdot  \tan\left(\dfrac{\pi}{3}-x\right)\cdot \tan\left(\dfrac{\pi}{3}+x\right)\\
		&=&\tan x\cdot \dfrac{\sqrt{3}-\tan x}{1+\sqrt{3}\tan x}\cdot \dfrac{\sqrt{3}+\tan  x}{1-\sqrt{3}\tan x}\\
		&=& \tan x\cdot  \dfrac{3-\tan^2x}{1-3\tan^2x}\\
		&=& \dfrac{3\tan  x-\tan^3x}{1-3\tan^2x}\\
		 &=&\tan 3x.
	\end{eqnarray*}
Vậy $n=3$.} 
\end{ex}

% nguồn: Ok_Toan11_TracNghiem_Lan2
\begin{ex}%[1D1H3-2]%[ma:7937b5dd]
	Cho $\tan \alpha=2$. Tính $\tan  \left(\alpha-\dfrac{\pi}{4}\right)$
	 \choice
		{$-\dfrac{1}{3}$} 
		 {$1$ }
		{$\dfrac{2}{3}$} 
		{\True $\dfrac{1}{3}$} 
	 \loigiai{
	Ta có  $\tan \left(\alpha-\dfrac{\pi}{4}\right)=\dfrac{\tan \alpha-\tan \dfrac{\pi}{4}}{1+\tan \alpha \tan \dfrac{\pi}{4}}=\dfrac{2-1}{1+2}=\dfrac{1}{3}$.}
\end{ex}

% nguồn: Ok_Toan11_TracNghiem_Lan2
\begin{ex}%[1D1H3-2]%[ma:eaf66125]
	Gọi $M = \dfrac{\sin(y - x)}{\sin x \cdot \sin y}$. Rút gọn biểu thức $M$, ta  được
	 \choice
		{$M = \tan x - \tan  y$}
		 {\True $M = \cot x - \cot y$}
		{$M = \cot  y - \cot x$}
		{$M = \dfrac{1}{\sin x} - \dfrac{1}{\sin y}$} 
	\loigiai{ 
		Ta có $M = \dfrac{\sin y\cdot\cos x - \cos y\cdot\sin x}{\sin x\cdot\sin y} = \dfrac{\sin  y\cdot\cos x}{\sin x\cdot\sin y} - \dfrac{\cos y\cdot\sin x}{\sin x\cdot\sin y} = \dfrac{\cos x}{\sin x} - \dfrac{\cos y}{\sin y} = \cot x - \cot y$.
	} 
\end{ex}

% nguồn: Ok_Toan11_TracNghiem_Lan2
\begin{ex}%[1D1H3-2]%[ma:75b1b598]
	Biểu thức $\dfrac{\sin (a+b)}{\sin (a-b)}$ bằng biểu thức nào  sau đây?
	\choice 
		{$\dfrac{\sin (a+b)}{\sin (a-b)}=\dfrac{\sin a+\sin  b}{\sin a-\sin b}$}
		{$\dfrac{\sin (a+b)}{\sin (a-b)}=\dfrac{\sin a-\sin b}{\sin  a+\sin b}$}
		 {\True $\dfrac{\sin (a+b)}{\sin (a-b)}=\dfrac{\tan a+\tan b}{\tan a-\tan b}$}
		{$\dfrac{\sin (a+b)}{\sin (a-b)}=\dfrac{\cot a+\cot b}{\cot  a-\cot b}$}
	\loigiai{ 
	Ta có $$ \dfrac{\sin (a+b)}{\sin (a-b)}=\dfrac{\sin a \cos b+\cos a \sin b}{\sin a \cos b-\cos a \sin b}=\dfrac{\tan a \cdot \cot b+1}{\tan  a \cdot \cot b-1}=\dfrac{\tan a \cdot \dfrac{1}{\tan b}+1}{\tan a \cdot \dfrac{1}{\tan b}-1}=\dfrac{\tan a+\tan b}{\tan a-\tan b}.$$}
\end{ex}

% nguồn: Ok_Toan11_TracNghiem_Lan2
\begin{ex}%[1D1H3-2]%[ma:4b737941]
	Rút gọn biểu thức: $\cos \left(120^{\circ}-\; x\right)+\cos \left(120^{\circ}+\; x\right)-\cos x$ ta được kết  quả là
	 \choice
		{$0$} 
		 {$-\cos x$}
		{\True $-2\cos  x$}
		 {$\sin x-\cos x$}
	 \loigiai{
		\begin{eqnarray*}
			&& \cos \left(120^{\circ}-\;  x\right)+\cos \left(120^{\circ}+\; x\right)-\cos x\\
			&=&-\dfrac{1}{2} \cos x+\dfrac{\sqrt{3}}{2} \sin x-\dfrac{1}{2} \cos x+\dfrac{\sqrt{3}}{2} \sin x-\cos x\\&=&-2\cos  x.
		\end{eqnarray*}
	} 
\end{ex}

% nguồn: Ok_Toan11_TracNghiem_Lan2
\begin{ex}%[1D1H3-2]%[ma:7b198077]
	Nếu biết $\sin a=\dfrac{8}{17},$ $\tan b=\dfrac{5}{12}$  và $a,$ $b$ đều là các góc nhọn và dương thì $\sin \left(a-b\right)$ bằng
	\choice 
		 {$\dfrac{20}{220}$}
		{$-\dfrac{20}{220}$} 
		{\True  $\dfrac{21}{221}$}
		{$\dfrac{22}{221}$} 
	 \loigiai{
		Ta có $a$, $b$ đều  là các góc nhọn và dương nên $\sin a,$ $ \sin b >0$.\\
		$\sin a=\dfrac{8}{17} \Rightarrow \cos  a=\sqrt{1-\dfrac{64}{289}}=\dfrac{15}{17}.$\\
		$\tan b=\dfrac{5}{12} \Rightarrow \cos b=\dfrac{1}{\sqrt{1+\dfrac{25}{144}}}=\dfrac{12}{13}. \\ \Rightarrow \sin b=\tan  b\cdot \cos b=\dfrac{5}{13}.$\\
		$\Rightarrow \sin \left(a-b\right)=\dfrac{8}{17}\cdot \dfrac{12}{13}-\dfrac{15}{17}\cdot  \dfrac{5}{13}=\dfrac{21}{221}.$
	} 
\end{ex}

% nguồn: Ok_Toan11_TracNghiem_Lan2
\begin{ex}%[1D1H3-2]%[ma:19a841e6]
	Cho $x$, $y$ là các  góc nhọn, $\cot x=\dfrac{4}{3}$, $\cot y=7$. Tổng $x+y$ bằng
	 \choice
		 {$\dfrac{\pi}{3}$}
		{\True  $\dfrac{\pi}{4}$}
		 {$\dfrac{\pi}{6}$}
		 {$\dfrac{2\pi}{3}$}
	\loigiai{Ta có $\cot x=\dfrac{4}{3} \Leftrightarrow \tan x=\dfrac{3}{4}$; $\cot  y=7 \Leftrightarrow \tan y=\dfrac{1}{7}$.\\
		Nên 
		$$\tan (x+y)=\dfrac{\tan x+\tan y}{1-\tan x \cdot \tan y}=1.$$ 
	 Suy ra $x+y=\dfrac{\pi}{4}$.}
\end{ex}

% nguồn: Ok_Toan11_TracNghiem_Lan2
\begin{ex}%[1D1H3-2]%[ma:a8d3d00f]
	Khẳng định nào đúng trong  các khẳng định sau?
	\choice 
		{$\sin a + \cos a = \sqrt{2}  \sin \left(a - \dfrac{\pi}{4}\right)$}
		{\True $\sin a + \cos a = \sqrt{2} \sin \left(a  + \dfrac{\pi}{4}\right)$}
		{$\sin a  + \cos a = -\sqrt{2} \sin \left(a - \dfrac{\pi}{4}\right)$}
		{$\sin a + \cos a = -\sqrt{2} \sin \left(a + \dfrac{\pi}{4}\right)$} 
	 \loigiai{
		Ta có $\sin a + \cos a = \sqrt{2} \left(\dfrac{1}{\sqrt{2}}\sin a  + \dfrac{1}{\sqrt{2}} \cos a\right)=\sqrt{2}\left(\cos\dfrac{\pi}{4}\sin a + \sin\dfrac{\pi}{4} \cos a\right)= \sqrt{2} \sin \left(a + \dfrac{\pi}{4}\right)$.
	} 
\end{ex}

% nguồn: Ok_TracNghiemLop11_Lan1
\begin{ex}%[1D1H3-2]%[ma:c17d7f11]
	Cho $\tan\alpha=3$.  Khi đó $\tan\left(\alpha-\dfrac{\pi}{4}\right)$ bằng
	 \choice
		{$-\dfrac{1}{2}$} 
		{$2$} 
		 {\True $-2$}
		 {$\dfrac{1}{2}$}
	 \loigiai{
		Ta có $\tan\left(\alpha-\dfrac{\pi}{4}\right)=\dfrac{\tan\alpha+\tan\dfrac{\pi}{4}}{1-\tan\alpha\cdot \tan\dfrac{\pi}{4}}=\dfrac{3+1}{1-3}=-2$. 
	} 
\end{ex}

% nguồn: Ok_TracNghiemLop11_Lan1
\begin{ex}%[1D1H3-2]%[ma:60b4c645]
	Rút gọn $M=\cos (a+b)  \cos (a-b)-\sin (a+b) \sin (a-b)$.
	\choice{$M=\sin 2a$} 
		{$M=\cos 4  a$}
		 {\True $M=\cos 2a$}
		 {$M=\sin 4a$}
	\loigiai{ 
		Áp dụng công thức cộng ta có $M=\cos (a+b) \cos  (a-b)-\sin (a+b) \sin (a-b)=\cos 2a$.
	} 
\end{ex}

% nguồn: Ok_TracNghiemLop11_Lan1
\begin{ex}%[1D1H3-2]%[ma:2beee978]
	Cho hai góc nhọn $a;b$ thoả  $\cos a=\dfrac{1}{3};\cos b=\dfrac{1}{4}.$ Tính giá trị của biểu thức $P=\cos \left({a+b}\right)\cdot\cos \left({a-b}\right).$
	 \choice
		{$-\dfrac{{113}}{{144}}$} 
		{$-\dfrac{{115}}{{144}}$} 
		 {$-\dfrac{{117}}{{144}}$}
		 {\True $-\dfrac{{119}}{{144}}$}
	\loigiai{ Ta có 
		\begin{eqnarray*}
			P&=&\cos \left({a+b}\right)\cdot\cos \left({a-b}\right)\\&=&\left({\cos a\cdot\cos b+\sin a\cdot\sin b}\right)\left({\cos a\cdot\cos b-\sin  a\cdot\sin b}\right)
			\\&=&\left({\cos  a\cdot\cos b}\right)^2-\left({\sin a\cdot\sin b}\right)^2\\&=&\cos ^2a\cdot\cos ^2b-\left({1-{\cos}^2a}\right)\cdot\left({1-{\cos}^2b}\right)
			 \\&=&\dfrac{1}{9}\cdot\dfrac{1}{{16}}-\left({1-\dfrac{1}{9}}\right)\cdot\left({1-\dfrac{1}{{16}}}\right)\\&=&-\dfrac{{119}}{{144}}.
		\end{eqnarray*}
	} 
\end{ex}

% nguồn: Ok_Toan11_TracNghiem_Lan2
\begin{ex}%[1D1V3-2]%[ma:b328c60c]
	Giá trị đúng của biểu thức $A=\dfrac{\tan 30^\circ+\tan 40^\circ+\tan 50^\circ+\tan 60^\circ}{\cos 20^\circ}$  bằng
	\choice 
		{$\dfrac{2}{\sqrt 3}$} 
		{$\dfrac{4}{\sqrt 3}$} 
		{$\dfrac{6}{\sqrt 3}$} 
		{\True $\dfrac{8}{\sqrt 3}$} 
	\loigiai{ 
		Ta có 
		\begin{align*}
			A&=\dfrac{\tan 30^\circ+\tan 40^\circ+\tan  50^\circ+\tan 60^\circ}{\cos 20^\circ}\\&=\dfrac{\dfrac{\sin 70^\circ}{\cos 30^\circ \cdot\cos 40^\circ}+\dfrac{\sin 110^\circ}{\cos 50^\circ \cdot\cos 60^\circ}}{\cos 20^\circ}\\
			&=\dfrac{1}{\cos 30^\circ \cdot\cos 40^\circ}+\dfrac{1}{\cos 50^\circ \cdot\cos  60^\circ}\\&=\dfrac{2}{\sqrt 3\cos 40^\circ}+\dfrac{2}{\cos 50^\circ}\\
			&=2\left(\dfrac{\cos 50^\circ+\sqrt 3\cos 40^\circ}{\sqrt 3\cos  40^\circ \cdot\cos 50^\circ}\right)\\&=2\left(\dfrac{\sin 40^\circ+\sqrt 3\cos 40^\circ}{\sqrt 3\cos 40^\circ \cdot\cos 50^\circ}\right)\\
			&=4\dfrac{\sin 100^\circ}{\dfrac{\sqrt 3}{2}\left(\cos 10^\circ+\cos 90^\circ\right)}\\&=\dfrac{8\cos 10^\circ}{\sqrt  3\cos 10^\circ}\\&=\dfrac{8}{\sqrt 3}.
		\end{align*}
	 }
\end{ex}

% nguồn: Ok_Toan11_TracNghiem_Lan2
\begin{ex}%[1D1V3-2]%[ma:b7484e99]
	Cho $M=3\sin x+4\cos  x$. Chọn khẳng định đúng.
	\choice 
		{\True  $-5\le M\le 5$}
		{$M  > 5$}
		 {$M\ge 5$}
		{$M\le  5$}
	\loigiai{ 
		$M=5\left(\dfrac{3}{5}\sin x+\dfrac{4}{5}{\cos x}\right)=5\sin(x+a)$ với $\cos  a=\dfrac{3}{5};\,\sin a=\dfrac{4}{5}$.\\
		Ta có $-1\le\sin(x+a)\le  1\Leftrightarrow-5\le 5\sin(x+a)\le 5$.
	} 
\end{ex}


%%%%%%%%%%%%%%%% ĐÚNG/SAI %%%%%%%%%%%%%%%%

% nguồn: Ok_dungsai_lop11_Lan1
\begin{ex}%[1D1N3-2]%[ma:fc35ed40]
  Mỗi khẳng định dưới đây đúng hay sai?
  \choiceTF
  {\True $\sin 7 a  \cdot \sin a+\cos 7 a \cdot \cos a=\cos 6a$}
  {\True $\sin 2 a \cdot \cos a-\cos 2a  \cdot \sin a=\sin a$}
  {\True  $\sin\left(a+\dfrac{\pi}{6}\right)=\dfrac{\sqrt{3}}{2} \sin a+\dfrac{1}{2}\cos a$}
  {\True Nếu $\tan \alpha =2$ thì giá  trị $\tan \left(\alpha - \dfrac{\pi}{4}\right)=\dfrac{1}{3}$}
 \loigiai{ 
  \begin{itemchoice}
   \itemch Ta có $\sin 7 a \cdot \sin a+\cos 7 a \cdot \cos a=\cos (7 a-a)=\cos 6 a$. 
   \itemch Ta có $\sin 2 a \cdot \cos a-\cos 2a \cdot \sin a  = \sin (2a -a)= \sin a$.
   \itemch Ta có $\sin\left(a+\dfrac{\pi}{6}\right)=\sin a \cos\dfrac{\pi}{6}+\cos a \sin \dfrac{\pi}{6}=\dfrac{\sqrt{3}}{2} \sin a+\dfrac{1}{2}\cos  a$.
   \itemch Ta có $\tan \left(\alpha - \dfrac{\pi}{4}\right)=\dfrac{\tan \alpha - \tan \dfrac{\pi}{4}}{1 + \tan \alpha \cdot \tan  \dfrac{\pi}{4}}=\dfrac{2-1}{1+2\cdot 1}=\dfrac{1}{3}$.
  \end{itemchoice}
  }
\end{ex}

% nguồn: Ok_dungsai_lop11_Lan1
\begin{ex}%[1D1N3-2]%[ma:b203c841]
 Mỗi khẳng  định dưới đây đúng hay sai?
 \choiceTF 
  {\True $\sin(a+b)+\sin(a-b)=2\sin a\cos b$} 
  {\True $\sin(a+b)-\sin(a-b)=2\cos a\sin b$} 
  {$\cos(a+b)+\cos(a-b)=2\sin  a\sin b$}
  {\True $\cos(a+b)-\cos(a-b)=-2\sin  a\cos b$}
  \loigiai{
  \begin{itemchoice}
   \itemch $\sin(a+b)+\sin(a-b)=(\sin a\cos b+\cos a\sin b)+(\sin a\cos b-\cos  a\sin b)=2\sin a\cos b$.
   \itemch $\sin(a+b)-\sin(a-b)=(\sin a\cos  b+\cos a\sin b)-(\sin a\cos b-\cos a\sin b)=2\cos a\sin b$.
   \itemch $\cos(a+b)+\sin(a-b)=(\cos a\cos b-\sin a\sin b)+(\cos a\cos b+\sin a\sin  b)=2\cos a\cos b$.
   \itemch $\cos(a+b)-\cos(a-b)=(\cos a\cos  b-\sin a\sin b)-(\cos a\cos b+\sin a\sin b)=-2\sin a\sin b$.
  \end{itemchoice}
  }
\end{ex}

% nguồn: Ok_dungsai_lop11_Lan2
\begin{ex}%[1D1N3-2]%[ma:cd942e50]
	Cho biết $\tan x =  \sqrt{2}$ và $0 < x < 90^\circ$.
	\choiceTF 
		{\True  $\cos x > 0$}
		 {\True $\cos x = \dfrac{\sqrt{3}}{3}$}
		{\True  $\sin x = \dfrac{\sqrt{6}}{3}$}
		 {$\cos \left(x - 30^\circ\right) = \dfrac{3 - \sqrt{6}}{6}$}
	\loigiai{ 
		\begin{itemchoice}
			\itemch Vì $0 < x < 90^\circ$  nên $\cos x > 0$.
			\itemch Ta có $\dfrac{1}{\cos^2 x} = 1 + \tan^2 x = 1 +  2 = 3 \Rightarrow \cos^2 x = \dfrac{1}{3} \Rightarrow \cos x = \dfrac{\sqrt{3}}{3}$.
			\itemch Mặt khác $\tan  x = \dfrac{\sin x}{\cos x} \Rightarrow \sin x = \tan x \cos x = \dfrac{\sqrt{6}}{3}$.
			\itemch Ta có $\cos(x - 30^\circ) = \cos x  \cos 30^\circ + \sin x \sin 30^\circ = \dfrac{\sqrt{3}}{3} \cdot \dfrac{\sqrt{3}}{2} + \dfrac{\sqrt{6}}{3} \cdot \dfrac{1}{2} = \dfrac{3+\sqrt{6}}{6}$.
		\end{itemchoice}
	} 
\end{ex}

% nguồn: Ok_dungsai_lop11_Lan2
\begin{ex}%[1D1N3-2]%[ma:76c6ae87]
	Biết  $\sin a = \dfrac{8}{17}$, $\tan b = \dfrac{5}{12}$ và $a$, $b$ là các góc nhọn.
	\choiceTF 
		{\True $\tan a  = \dfrac{8}{15}$}
		{\True $\sin (a-b)  = \dfrac{21}{221}$}
		{$\cos  (a+b) = \dfrac{14}{22}$}
		{$\tan (a+b) =  \dfrac{17}{14}$}
	 \loigiai{
		\begin{itemchoice}
			\itemch Vì  $a,b$ là các góc nhọn nên $\cos a > 0, \cos b > 0$.\\
			Ta có $\cos a = \sqrt{1-\sin^2 a} = \dfrac{15}{17} \Rightarrow  \tan a = \dfrac{\sin a}{\cos a} = \dfrac{8}{15}$.
			\itemch Ta có $\cos b =  \sqrt{\dfrac{1}{1+\tan^2 b}} = \dfrac{12}{13} \Rightarrow \sin b = \cos b \tan b = \dfrac{5}{13}$.\\
			Khi đó $\sin(a-b) = \sin a \cos b - \cos a \sin b = \dfrac{8}{17} \cdot  \dfrac{12}{13} - \dfrac{15}{17} \cdot \dfrac{5}{13} = \dfrac{21}{221}$.
			\itemch Ta có $\cos(a+b) = \cos a \cos b - \sin a \sin b = \dfrac{15}{17} \cdot \dfrac{12}{13}  - \dfrac{8}{17} \cdot \dfrac{5}{13} = \dfrac{140}{221}$.
			\itemch Ta  có $\tan(a+b) = \dfrac{\tan a + \tan b}{1 - \tan a \tan b} = \dfrac{\dfrac{8}{15} + \dfrac{5}{12}}{1 - \dfrac{8}{15} \cdot \dfrac{5}{12}} = \dfrac{171}{140}$.
		\end{itemchoice}
	} 
\end{ex}

% nguồn: Ok_dungsai_lop11_Lan2
\begin{ex}%[1D1N3-2]%[ma:92f1c752]
	Cho biết $\sin  \alpha = -\dfrac{12}{13}$, $\dfrac{3\pi}{2} < \alpha < 2\pi$.
	 \choiceTF
		 {\True $\cos \alpha > 0$}
		 {\True $\cos \alpha = \dfrac{5}{13}$}
		{\True $\tan \alpha  = -\dfrac{12}{5}$}
		{$\cos  \left(\dfrac{\pi}{3} - \alpha\right) = \dfrac{5-\sqrt{3}}{26}$}
	\loigiai{ 
		\begin{itemchoice}
			\itemch Ta có $\cos\left(\dfrac{\pi}{3} - \alpha\right) = \cos\dfrac{\pi}{3}\cos\alpha + \sin\dfrac{\pi}{3}\sin\alpha= \dfrac{1}{2}\cos\alpha + \dfrac{\sqrt{3}}{2} \cdot  \left(-\dfrac{12}{13}\right) = \dfrac{1}{2}\cos\alpha - \dfrac{6\sqrt{3}}{13}$.\\
			Do  $\dfrac{3\pi}{2} < \alpha < 2\pi \Rightarrow \cos\alpha > 0$.
			\itemch Ta có $\cos\alpha = \sqrt{1-\sin^2 \alpha}  = \sqrt{1-\left(-\dfrac{12}{13}\right)^2} = \sqrt{1-\dfrac{144}{169}} = \sqrt{\dfrac{25}{169}} = \dfrac{5}{13}$.
			\itemch Ta có $\tan\alpha = \dfrac{\sin\alpha}{\cos\alpha} = \dfrac{-\dfrac{12}{13}}{\dfrac{5}{13}} =  -\dfrac{12}{5}$.
			\itemch Ta có $\cos\left(\dfrac{\pi}{3}  - \alpha\right) = \dfrac{1}{2}\cos\alpha + \dfrac{\sqrt{3}}{2}\sin\alpha = \dfrac{1}{2} \cdot \dfrac{5}{13} + \dfrac{\sqrt{3}}{2} \cdot \left(-\dfrac{12}{13}\right) = \dfrac{5}{26} - \dfrac{12\sqrt{3}}{26} = \dfrac{5-12\sqrt{3}}{26}$.
		\end{itemchoice}
	} 
\end{ex}

% nguồn: Ok_dungsai_lop11_Lan2
\begin{ex}%[1D1N3-2]%[ma:9dc0b8ec]
	Cho biết $\sin x = \dfrac{1}{\sqrt{3}}$ và $0 < x <  \dfrac{\pi}{2}$.
	\choiceTF 
		 {\True $\cos x > 0$}
		{\True  $\cos x = \dfrac{\sqrt{6}}{3}$}
		 {$\tan x = \dfrac{\sqrt{3}}{3}$}
		 {$\cos \left(x + \dfrac{\pi}{3}\right) = \dfrac{\sqrt{6}-3}{8}$}
	\loigiai{ 
		\begin{itemchoice}
			 \itemch Vì $0 < x < \dfrac{\pi}{2}$ nên $\cos x > 0$.
			\itemch Ta có $\sin x = \dfrac{1}{\sqrt{3}} \Rightarrow \cos x = \sqrt{1-\sin^2 x}  = \sqrt{1-\dfrac{1}{3}} = \dfrac{\sqrt{6}}{3}$.
			\itemch  Ta có $\tan x = \dfrac{\sin x}{\cos x} = \dfrac{\sqrt{2}}{2}$.
			\itemch Ta có $\cos\left(x + \dfrac{\pi}{3}\right) = \cos x \cos \dfrac{\pi}{3} - \sin x \sin \dfrac{\pi}{3} = \dfrac{\sqrt{6}}{3} \cdot \dfrac{1}{2}  - \dfrac{1}{\sqrt{3}} \cdot \dfrac{\sqrt{3}}{2} = \dfrac{\sqrt{6}-3}{6}$.
		\end{itemchoice}
	} 
\end{ex}

% nguồn: Ok_dungsai_lop11_Lan2
\begin{ex}%[1D1N3-2]%[ma:610c1281]
	Cho biết $\sin \alpha = \dfrac{3}{5}$, $\dfrac{\pi}{2} < \alpha  < \pi$.
	\choiceTF 
		{\True $\cos \alpha  < 0$}
		{\True $\cos  \alpha = -\dfrac{4}{5}$}
		 {$\tan \alpha = \dfrac{3}{4}$}
		{$\tan \left(\alpha + \dfrac{\pi}{3}\right) = \dfrac{48  - \sqrt{3}}{11}$}
	\loigiai{ 
		\begin{itemchoice}
			\itemch Vì $\dfrac{\pi}{2} <  \alpha < \pi \Rightarrow \cos\alpha < 0$.
			\itemch  Ta có $\cos\alpha = -\sqrt{1-\sin^2 \alpha} = -\dfrac{4}{5}$.
			\itemch Ta  có $\tan\alpha = \dfrac{\sin\alpha}{\cos\alpha} = -\dfrac{3}{4}$.
			\itemch Ta có $\tan\left(\alpha + \dfrac{\pi}{3}\right) = \dfrac{\tan\alpha + \tan \dfrac{\pi}{3}}{1 -  \tan\alpha \tan \dfrac{\pi}{3}} = \dfrac{\tan\alpha + \sqrt{3}}{1 - \sqrt{3} \tan\alpha}$.\\
			Suy ra  $\tan\left(\alpha + \dfrac{\pi}{3}\right) = \dfrac{-\dfrac{3}{4} + \sqrt{3}}{1 - \sqrt{3}\left(-\dfrac{3}{4}\right)} = \dfrac{\dfrac{-3+4\sqrt{3}}{4}}{\dfrac{4+3\sqrt{3}}{4}} = \dfrac{-3+4\sqrt{3}}{4+3\sqrt{3}} =
			\dfrac{48 - 25\sqrt{3}}{11}$. 
		\end{itemchoice}
	 }
\end{ex}

% nguồn: Ok_dungsai_lop11_Lan2
\begin{ex}%[1D1H3-2]%[ma:a097b7b5]
	Cho biết $\cos x = -\dfrac{12}{13}$ và $\pi < x <  \dfrac{3\pi}{2}$.
	 \choiceTF
		{$\sin x > 0$} 
		{\True $\sin x = -\dfrac{5}{13}$} 
		{$\cot  x = \dfrac{5}{12}$}
		{\True $\sin  \left(\dfrac{\pi}{3} - x\right) = \dfrac{5-12\sqrt{3}}{26}$}
	 \loigiai{
		\begin{itemchoice}
			\itemch Vì $\pi < x < \dfrac{3\pi}{2}$ nên  $\sin x < 0$.
			\itemch Ta có $\cos x = -\dfrac{12}{13} \Rightarrow \sin x  = -\sqrt{1-\cos^2 x} = -\sqrt{1-\left(-\dfrac{12}{13}\right)^2} = -\dfrac{5}{13}$.
			\itemch Ta  có $\cot x = \dfrac{\cos x}{\sin x} = \dfrac{12}{5}$.
			\itemch Ta có $\sin\left(\dfrac{\pi}{3} -  x\right) = \sin\dfrac{\pi}{3}\cos x - \cos\dfrac{\pi}{3}\sin x = \dfrac{\sqrt{3}}{2} \cdot \left(-\dfrac{12}{13}\right) - \dfrac{1}{2} \cdot \left(-\dfrac{5}{13}\right)= \dfrac{5-12\sqrt{3}}{26}$.
		\end{itemchoice}
	} 
\end{ex}

% nguồn: Ok_dungsai_lop11_Lan1
\begin{ex}%[1D1V3-2]%[ma:dca649de]
 Biết $\sin a=\dfrac{8}{17}, \tan b=\dfrac{5}{12}$ và $a$, $b$ là các góc nhọn. Khi đó 
  \choiceTF
  {\True  $\tan a=\dfrac{8}{15}$}
   {\True $\sin (a-b)=\dfrac{21}{221}$}
  {$\cos  (a+b)=\dfrac{14}{22}$}
  {$\tan (a+b)=\dfrac{17}{14}$} 
 \loigiai{ 
  \begin{itemchoice}
   \itemch {\bf  Đúng.}\\
   Vì $a$, $b$ là các góc nhọn nên $\cos  a > 0,\cos b > 0$.\\
   Ta có $\cos a=\sqrt{1-\sin^2 a}=\dfrac{15}{17} \Rightarrow \tan a=\dfrac{\sin  a}{\cos a}=\dfrac{8}{15}$.
   \itemch {\bf  Đúng.}\\
   $\cos b=\sqrt{\dfrac{1}{1+\tan^2 b}}=\dfrac{12}{13} \Rightarrow \sin b=\cos b\tan b=\dfrac{5}{13}$.\\ 
   Khi đó $\sin (a-b)=\sin a\cos b-\cos a\sin b=\dfrac{8}{17} \cdot  \dfrac{12}{13}-\dfrac{15}{17} \cdot \dfrac{5}{13}=\dfrac{21}{221}$.
   \itemch  {\bf Sai.}\\
   $\cos (a+b)=\cos a\cos  b-\sin a\sin b=\dfrac{15}{17} \cdot \dfrac{12}{13}-\dfrac{8}{17} \cdot \dfrac{5}{13}=\dfrac{140}{221}$.
   \itemch  {\bf Sai.}\\
   $\tan  (a+b)=\dfrac{\tan a+\tan b}{1-\tan a\tan b}=\dfrac{\dfrac{8}{15}+\dfrac{5}{12}}{1-\dfrac{8}{15} \cdot \dfrac{5}{12}}=\dfrac{171}{140}$.
  \end{itemchoice}
 } 
\end{ex}


%%%%%%%%%%%%%%%% TRẢ LỜI NGẮN %%%%%%%%%%%%%%%%

% nguồn: Ok_TraLoiNgan_Lop11_Lan2
\begin{ex}%[1D1H3-2]%[ma:1dca6c1f]
	Cho $\alpha-\beta=\dfrac{\pi}{3}$ . Tính giá trị của biểu thức  $A=(\cos\alpha+\cos\beta)^2+(\sin\alpha+\sin\beta)^2$.
	 \shortans{3}
	\loigiai{ 
		\allowdisplaybreaks 
		\begin{eqnarray*}	A&=&\cos^2\alpha+\cos^2\beta+2\cos\alpha\cos\beta+\sin^2\alpha+\sin^2\beta+2\sin\alpha\sin\beta\\
			 &=&2+2\cos\alpha\cos\beta+2\sin\alpha\sin\beta=2+2\cos (\alpha-\beta)=2+2\cos\dfrac{\pi}{3}=3.
		\end{eqnarray*}
	 }
\end{ex}

% nguồn: Ok_TraLoiNgan_Lop11_Lan2
\begin{ex}%[1D1H3-2]%[ma:bdb35e19]
	Cho hai góc nhọn $a$ và $b$. Biết $\cos a=\dfrac{1}{3}$, $\cos b=\dfrac{1}{4}$. Tính giá trị của biểu thức $P=\cos (a+b)\cdot\cos (a-b)$ (Kết quả làm tròn  đến hàng phần chục).
	\shortans{-0,8} 
	\loigiai{ 
		Ta có  $\sin ^2a=1-\cos ^2a=\dfrac{8}{9}$, ${\sin ^2}b=1-\cos ^2b=\dfrac{15}{16}$.\\
		\allowdisplaybreaks 
		\begin{eqnarray*}
			P&=&(\cos a\cos b-\sin  a\sin b)(\cos a\cos b+\sin a\sin b)\\
			&=&(\cos a\cos b)^2-(\sin a\sin b)^2=\cos^2a{\cos^2}b-\sin^2a{\sin^2}b\\ 
			&=&\dfrac{1}{9}\cdot\dfrac{1}{16}-\dfrac{8}{9}\cdot\dfrac{15}{16}=-\dfrac{119}{144}\approx  -0{,}8.
		\end{eqnarray*}
	 }
\end{ex}

% nguồn: Ok_TraLoiNgan_Lop11_Lan2
\begin{ex}%[1D1H3-2]%[ma:b213a1c4]
	Cho các góc $\alpha ,\beta  $ thỏa mãn $\dfrac{\pi}{2}<\alpha$, $\beta <\pi$, $\sin\alpha=\dfrac{1}{3}$, $\cos\beta=-\dfrac{2}{3}$.\\
	Tính $\sin (\alpha+\beta)$ (Kết quả  làm tròn đến hàng phần chục).
	 \shortans{-0,9}
	 \loigiai{
		Do  $\dfrac{\pi}{2}<\alpha ,\beta <\pi\Rightarrow\cos\alpha < 0,\;\sin\beta>0$.\\
		 $\cos\alpha=-\sqrt{1-\sin^2\alpha}=-\sqrt{1-\dfrac{1}{9}}=-\dfrac{2\sqrt 2}{3}$.\\
		$\sin\beta=\sqrt{1-\cos^2\beta}=\sqrt{1-\dfrac{4}{9}}=\dfrac{\sqrt 5}{3}$.\\ 
		$\sin (\alpha+\beta)=\sin\alpha\cdot\cos\beta+\cos\alpha\cdot\sin\beta=\dfrac{1}{3}\cdot\left(-\dfrac{2}{3}\right)+\left(-\dfrac{2\sqrt  2}{3}\right)\cdot\dfrac{\sqrt 5}{3}=-\dfrac{2+2\sqrt{10}}{9}$.\\
		Vậy  $\sin (\alpha+\beta)=-\dfrac{2+2\sqrt{10}}{9}\approx -0{,}9$.
	 }
\end{ex}

% nguồn: Ok_TraLoiNgan_Lop11_Lan2
\begin{ex}%[1D1H3-2]%[ma:e3601b40]
	Cho hai  góc nhọn $a$ và $b$ với $\tan a=\dfrac{1}{7}$ và $\tan b=\dfrac{3}{4}$. Tính $a+b$.
	\shortans{45} 
	\loigiai{ 
		Ta có $\tan (a+b)=\dfrac{\tan a+\tan  b}{1-\tan a\tan b}=\dfrac{\dfrac{1}{7}+\dfrac{3}{4}}{1-\dfrac{1}{7}\cdot\dfrac{3}{4}}=1$, suy ra $a+b=45^\circ$.
	} 
\end{ex}

% nguồn: Ok_TraLoiNgan_Lop11_Lan2
\begin{ex}%[1D1H3-2]%[ma:5749a7d3]
	Cho $\alpha-\beta=\dfrac{\pi}{3}$. Tính giá trị của biểu thức  sau (Kết quả làm tròn đến hàng phần trăm). $$B=(\cos\alpha+\sin\beta)^2+(\cos\beta-\sin\alpha)^2$$
	 \shortans{0,27}
	 \loigiai{
		\allowdisplaybreaks 
		\begin{eqnarray*}	B&=&\cos^2\alpha+\sin^2\beta+2\cos\alpha\sin\beta+\cos^2\beta+\sin^2\alpha-2\sin\alpha\cos\beta\\
			&=&  2+2\sin\beta\cos\alpha-2\cos\beta\sin\alpha\\
			 &=& 2+2\sin(\beta-\alpha)=2+2\sin\left(-\dfrac{\pi}{3}\right)=2-\sqrt3\approx 0{,}27.
		\end{eqnarray*}
	} 
\end{ex}

% nguồn: Ok_TraLoiNgan_Lop11_Lan2
\begin{ex}%[1D1V3-2]%[ma:710a4f21]
	Cho tam giác $ABC$ vuông tại $A$ có cạnh $AB=6,AC=8$ . Điểm $E$ thuộc đoạn $AC$ sao cho $\widehat{CBE}=30^\circ$, điểm $D$ thuộc tia đối của tia $BA$ sao cho $\widehat{BCD}=30^\circ$. Tính độ dài cạnh $AD$ (Kết quả  trong bài được làm tròn đến hàng phần chục).
	\begin{center}
		\begin{tikzpicture}[scale=0.6]
				 \node at (0,-0.4) {A};
				\node at  (6,-0.4) {B};
				\node  at (0,8.4) {C};
				\node at (18.7,-0.4)  {D};
				\node at (-0.4,2.56) {E}; 
				 \draw (0,0) -- (0,8);
				\draw (0,0) -- (18.7,0); 
				\draw (0,2.56) -- (6,0); 
				\draw  (0,8) -- (6,0);
				 \draw (18.7,0) -- (0,8);
				 \draw[<->] (0,-0.8) -- (6,-0.8);
				\node at (3,-1.3) {6}; 
				\draw[<->] (-0.8,0)  -- (-0.8,8);
				\node at  (-1.1,4) {8};
			\end{tikzpicture}
	\end{center}
	\shortans{18,7} 
	 \loigiai{
		Ta có $\tan  \widehat{ACD}=\tan \left(\widehat{ACB}+30^{\circ}\right)=\dfrac{\tan \widehat{ACB}+\tan 30^{\circ}}{1-\tan \widehat{ACB} \tan 30^{\circ}}=\dfrac{\dfrac{0}{8}+\dfrac{1}{\sqrt{3}}}{1-\dfrac{6}{8} \cdot \dfrac{1}{\sqrt{3}}}=\dfrac{48+25 \sqrt{3}}{39}$.\\
		 Tam giác $ACD$ vuông tại $A$ có $AD=AC \cdot\tan \widehat{ACD}=8 \cdot \dfrac{48+25 \sqrt{3}}{39} \approx 18,7$.
	} 
\end{ex}

% nguồn: Ok_TraLoiNgan_Lop11_Lan2
\begin{ex}%[1D1V3-2]%[ma:eac9a039]
	Cho $\sin \alpha=\dfrac{3}{5}, \cos \beta=\dfrac{12}{13}$ và $0^{\circ} < \alpha, \beta < 90^{\circ}$. Biết $P=\sin \left(\alpha+\beta \right)$ và $Q=\cos \left(\alpha-\beta \right)$. Khi đó hãy tính giá trị của $P+Q$ (kết  quả làm tròn đến hàng phần trăm).\shortans{1,83}
	 \loigiai{
		Vì $0^{\circ} < \alpha  < 90^{\circ}$ nên $\cos \alpha > 0$ $\Rightarrow \cos \alpha=\sqrt{1-\sin ^2\alpha}=\sqrt{1-\left(\dfrac{3}{5} \right)^2}=\dfrac{4}{5}$.\\
		Vì $0^{\circ} < \beta < 90^{\circ}$ nên $\sin \beta > 0$ $\Rightarrow  \sin \beta=\sqrt{1-\cos ^2\beta}=\sqrt{1-\left(\dfrac{12}{13} \right)^2}=\dfrac{5}{13}$.\\
		Khi đó: $P=\sin \left(\alpha+\beta  \right)=\sin \alpha \cos \beta+\cos \alpha \sin \beta=\dfrac{3}{5} \cdot \dfrac{12}{13}+\dfrac{4}{5} \cdot \dfrac{5}{13}=\dfrac{56}{65}.$\\
		$Q=\cos \left(\alpha-\beta \right)=\cos \alpha \cos \beta+\sin \alpha \sin \beta=\dfrac{4}{5} \cdot \dfrac{12}{13}+\dfrac{3}{5} \cdot \dfrac{5}{13}=\dfrac{63}{65}.$\\ 
		Suy ra giá  trị biểu thức $P+Q=\dfrac{56}{65}+\dfrac{63}{65}=\dfrac{119}{65} \approx 1{,}83$.
	 }
\end{ex}

% nguồn: Ok_TraLoiNgan_Lop11_Lan2
\begin{ex}%[1D1V3-2]%[ma:c5c69a40]
	Cho tam giác $ABC$ vuông tại $A$ có cạnh $AB=6,AC=8$ . Điểm $E$ thuộc đoạn $AC$ sao cho $\widehat{CBE}=30^\circ$, điểm $D$ thuộc tia đối của tia $BA$ sao cho $\widehat{BCD}=30^\circ$. Tính độ dài cạnh $AE$ (Kết quả trong bài được  làm tròn đến hàng phần trăm).
	\begin{center}
	\end{center}
	 \shortans{2,56}
	\loigiai{ 
		Ta có $\tan\widehat{ABE}=\tan\left(\widehat{ABC}-30^\circ\right)=\dfrac{\tan\widehat{ABC}-\tan{30^\circ}}{1+\tan\widehat{ABC}\tan{30^\circ}}=\dfrac{\dfrac{8}{6}-\dfrac{1}{\sqrt 3}}{1+\dfrac{8}{6}\cdot\dfrac{1}{\sqrt 3}}=\dfrac{48-25\sqrt  3}{11}$.\\
		Tam giác $ABE$ vuông tại $A$ có $AE=AB\cdot\tan\widehat{ABE}=6\cdot\dfrac{48-25\sqrt  3}{11}\approx 2,56$.
	} 
\end{ex}


%%%%%%%%%%%%%%%% TỰ LUẬN %%%%%%%%%%%%%%%%

% nguồn: Ok_TuanLuan_Toan11_Lan2
\begin{ex}%[1D1H3-2]%[ma:cdd444a0]
	 Tính:
	\begin{listEX}
		 \item $\cos\left(a+\dfrac{\pi}{6}\right)$ biết $\sin a = \dfrac{1}{\sqrt{3}}$ và $\dfrac{\pi}{2} < a < \pi$.
		\item $\tan\left(a-\dfrac{\pi}{4}\right)$ biết $\cos a = -\dfrac{1}{3}$ và  $\pi < a < \dfrac{3\pi}{2}$.
	\end{listEX}
	 \loigiai{
		\begin{listEX}
			\item Vì $\dfrac{\pi}{2}  < a < \pi$ nên $\cos a < 0$.\\
			Mặt khác, từ $\sin^2 a + \cos^2 a = 1$ suy ra $\cos a  = -\sqrt{1 - \sin^2 a} = -\sqrt{1 - \left(\dfrac{1}{\sqrt{3}}\right)^2} = -\dfrac{\sqrt{6}}{3}$.\\
			Ta có $\cos\left(a+\dfrac{\pi}{6}\right) = \cos a  \cos \dfrac{\pi}{6} - \sin a \sin \dfrac{\pi}{6} = \left(-\dfrac{\sqrt{6}}{3}\right) \cdot \dfrac{\sqrt{3}}{2} - \dfrac{1}{\sqrt{3}} \cdot \dfrac{1}{2} = \dfrac{-3\sqrt{2}-\sqrt{3}}{6}$.
			\item Vì $\pi < a <  \dfrac{3\pi}{2}$ nên $\sin a < 0$, do đó $\tan a = \dfrac{\sin a}{\cos a} > 0$.\\
			Mặt khác $1 + \tan^2 a = \dfrac{1}{\cos^2  a}$ suy ra $\tan a = \sqrt{\dfrac{1}{\cos^2 a} - 1} = \sqrt{\dfrac{1}{\left(-\dfrac{1}{3}\right)^2} - 1} = 2\sqrt{2}$.\\
			Ta có $\tan\left(a-\dfrac{\pi}{4}\right) = \dfrac{\tan a - \tan \dfrac{\pi}{4}}{1 + \tan a \tan \dfrac{\pi}{4}} = \dfrac{2\sqrt{2} - 1}{1 + 2\sqrt{2}  \cdot 1} = \dfrac{9 - 4\sqrt{2}}{7}$.
		\end{listEX}
	 }
\end{ex}
